\documentclass[12pt,a4paper]{article}
\usepackage{amsmath,amssymb}
\usepackage{hyperref}
\hypersetup{colorlinks=true,linkcolor=blue,urlcolor=blue}
\usepackage{geometry}
\geometry{left=2.5cm,right=2.5cm,top=2.5cm,bottom=2.5cm}

\title{\textbf{L²/t Page Replacement: A Unified Frequency-Age Policy for Memory Management}}
\author{Stanislav Usychenko\\
\small Independent Researcher\\
\small Dnipro, Ukraine\\
\href{mailto:stanislavusicenko@gmail.com}{stanislavusicenko@gmail.com}
}
\date{May 2026}

\begin{document}

\maketitle

\begin{abstract}
We propose a novel page replacement policy based on the viability score \(v = L^2 / (t + 1)\), where \(L\) is the access frequency and \(t\) is the age (time since last access) of a memory page. Unlike LRU (recency only) or LFU (frequency only), our approach combines both metrics nonlinearly. The quadratic term strongly protects frequently accessed pages, while the linear age ensures that old pages are eventually evicted. We evaluate the policy on the MSR Cambridge web trace and compare it against LRU and LFU. Results show a 5–7\% higher hit ratio than LRU and 1–2\% higher than LFU. The source code and a Linux kernel patch are available at \url{https://github.com/Apuoxo/l2t-page-replacement}.
\end{abstract}

\section{Introduction}
Page replacement is a classic problem in operating systems. Two dominant algorithms are LRU (evicts least recently used page) and LFU (evicts least frequently used page). Both have weaknesses: LRU ignores frequency, while LFU suffers from ``stale'' pages that were once popular but are no longer needed. We introduce a new policy based on a simple nonlinear combination:

\[
v = \frac{L^2}{t + 1}
\]

where \(L\) is the page access count (frequency) and \(t\) is the time (in jiffies) since the last access. When a page must be evicted, we choose the one with the smallest \(v\).

\section{The L²/t Principle}
\subsection{Motivation}
Consider two pages:
\begin{itemize}
    \item Page A: accessed 10 times, last access 1 second ago.
    \item Page B: accessed 1 time, last access 5 seconds ago.
\end{itemize}
LRU would evict Page A (older), which is wrong. LFU would evict Page B, which is better.
Now consider:
\begin{itemize}
    \item Page C: accessed 10 times, last access 100 seconds ago.
    \item Page D: accessed 5 times, last access 1 second ago.
\end{itemize}
LFU would keep Page C (higher frequency), but Page C is very old. L²/t gives:
\[
v_C = 100/101 \approx 0.99, \quad v_D = 25/2 = 12.5
\]
Thus Page C is evicted first, which is desirable.

\subsection{Why Square?}
Linear ratios \(L/t\) do not sufficiently protect high-frequency pages. The quadratic term creates a super-linear advantage for pages with \(L \ge 5\), matching the intuition that a page accessed 10 times is about four times as important as one accessed 5 times.

\section{Simulation Results}
We used the MSR Cambridge Traces (web1 trace, first 500,000 accesses). Cache size = 1000 pages.

\begin{table}[h]
\centering
\caption{Hit ratio comparison}
\begin{tabular}{|c|c|}
\hline
\textbf{Policy} & \textbf{Hit Ratio} \\
\hline
LRU  & 0.8124 \\
LFU  & 0.8451 \\
L²/t & 0.8612 \\
\hline
\end{tabular}
\end{table}

L²/t outperforms LRU by 4.88 percentage points and LFU by 1.61 percentage points.

\section{Linux Kernel Patch}
A proof-of-concept patch is available in the GitHub repository. The core eviction function scans 16 candidate pages and selects the one with smallest \(v\).

\section{Conclusion}
We have introduced L²/t page replacement, a simple policy that outperforms LRU and LFU on real traces. The source code, simulation scripts, and a kernel patch are available at \url{https://github.com/Apuoxo/l2t-page-replacement}.

\end{document}